\section{Conclusion}\label{conc}

This project set out to remove the experience-level friction from a capable but ageing assessment tool, and to do so on the live system rather than by replacing it. Each of the five objectives was met. A reproducible Docker environment now rebuilds Moodle, CodeRunner and its dependencies with a single command and runs both test harnesses, providing the regression safety net the rest of the work depended on. The author page gained flexible test-case management and a corrected Precheck display, together with the visual and accessibility improvements the analysis called for. The student page gained a switchable, resizable layout and reclaimable working space. Task-focused documentation, searchable and reachable from the editing form, was served from inside the plugin. The delivered changes were then evaluated with real students and teachers, whose feedback both validated the work and generated further requirements.

The evidence supports the project's central claim: the friction that remains in a mature tool of this kind is a software problem, and disciplined user-experience engineering resolves it. Every delivered feature won majority approval from both audiences, none was rejected, the accessibility score improved, and an independent regression suite showed that nothing already working was broken. Three requirements were deferred, each recorded as future work rather than abandoned. The clearest recommendation that follows is to offer the changes upstream to the CodeRunner project, so that a far larger community of authors and students benefits from them and subjects them to wider review, and to continue the evidence-led approach for the feedback and hint features that users asked for most.

